\documentclass{uflamon}          % classe base para a monografia

%==============================================================================
% Utilizacao de pacotes
\usepackage[T1]{fontenc}         % usa fontes postscript com acentos
\usepackage[brazil]{babel}       % hifenização e títulos em português do Brasil
\usepackage[utf8]{inputenc}     % permite edição direta com acentos
\usepackage{amsmath}             % pacote da AMS para Matemática Avançada
\usepackage{amssymb}             % símbolos extras da AMS
\usepackage{latexsym}            % símbolos extras do LaTeX
\usepackage{graphicx}            % para inserção de gráficos
\usepackage{listings}            % para inserção de código
\usepackage{fancyvrb}            % para inserção de saídas de comandos
%\usepackage{enumerate}           % para personalizar lista enumeradas 
											%(incluso na classe)
\usepackage{longtable}           % para tambelas muito grandes NOVO!!!!

\usepackage{colortbl} % cores em tabelas
\newcolumntype{Z}{|>{\columncolor[gray]{0.9}}l|} %cor cinza em células
%\usepackage{array} % já incluso na classe
\newcolumntype{L}[1]{>{\raggedright\let\newline\\\arraybackslash\hspace{0pt}}m{#1}}
\newcolumntype{C}[1]{>{\centering\let\newline\\\arraybackslash\hspace{0pt}}m{#1}}
\newcolumntype{R}[1]{>{\raggedleft\let\newline\\\arraybackslash\hspace{0pt}}m{#1}}
\usepackage{multirow} % para juntar duas linhas em uma só
\usepackage{multicol} % para uso de várias colunas
%\usepackage{subfig}

% cores para os links cruzados
\usepackage{color}
\definecolor{rltred}{rgb}{0.2,0,0}
\definecolor{rltgreen}{rgb}{0,0.2,0}
\definecolor{rltblue}{rgb}{0,0,0.2}

\usepackage[colorlinks=true,
            urlcolor=rltblue,       % \href{...}{...} external (URL)
            filecolor=rltgreen,     % \href{...} local file
            linkcolor=black,       % \ref{...} and \pageref{...}
            citecolor=black,
            pdftitle={Determinação da Perda de Carga por Acessórios},
          pdfauthor={Marcus Bruno Fernandes Silva},
          pdfsubject={Determinação da Perda de Carga por Acessórios},
          pdfkeywords={perda de carga}%
]{hyperref} % para referência cruzadas
%\usepackage{hyperref}            % para referência cruzadas
\usepackage{subfigure}           % figuras dentro de figuras
\usepackage{caption}            % remodelando o formato dos títulos de 
                                 % tabelas e figuras

% configuração padrão do listings   
\lstset{%
   language=Java,
   extendedchars=true,
   tabsize=3,
   basicstyle=\footnotesize\ttfamily,
   stringstyle=\em,
   showstringspaces=false 
}

%abnt-emphasize=bf coloca o título das bibliografias em negrito
%abnt-thesis-year=both
\usepackage[alf,num,abnt-etal-cite=3,abnt-etal-list=3,abnt-url-package=url,abnt-emphasize=bf]{abntex2cite}

% redefinindo formatação de títulos de tabelas e figuras

%==============================================================================
\usepackage{microtype} 			% para melhorias de justificação
\usepackage{float}
\usepackage{titling}
\usepackage[version=4]{mhchem}
\usepackage{array}
\usepackage{booktabs}
% \usepackage{nameref} % Referencia capitulos por nome
%\usepackage[authoryear]{natbib}

% Especificando hifenizações que por ventura LaTeX não saiba fazer
% Por padrão 99,9% dos termos em português devem ser hifenizados corretamente.
\hyphenation{hardware software Li-nux am-bien-te diag-nos-ti-car coor-de-na-ção 
FAE-PE Recovery TelEduc Williams UFLA Sérgio Renata Reynolds disciplina}

%==============================================================================
% Dados da monografia, capa: autor, titulo, banca, etc... - SUBSTITUA DE ACORDO
%==============================================================================
\author{Felipe Alexander Julio \\
       Gabriela Nunes \\
       Marcus Bruno Fernandes Silva \\
       Petrus Damasceno}
\subtitle{Exemplo para os Usuários}
%\author{Joao}
\title{Determinação da Perda de Carga por Acessórios}
\engtitle{Use of Uflamon Class}
%\engsubtitle{Sample for Users}
%\edicao{3$^a$ edição revista, atualizada e ampliada}
\date{2017}
\tipo{Projeto entregue e apresentado à Professora Renata de Aquino B. L. Corrêa como requisito da disciplina GNE334.}
\orientador{Professora Renata de Aquino B. L. Corrêa}
\newcommand{\orient}{Professora Renata de Aquino B. L. Corrêa}
\newcommand{\otipo}{Projeto entregue e apresentado à Professora Renata de Aquino B. L. Corrêa como requisito da disciplina GNE334.}
\newcommand{\olocal}{Lavras -- MG}
\local{Lavras -- MG}

%##################################################

\newcommand \fazercapa{
\begin{titlepage}
    \centering

    \includegraphics[scale=0.45]{logoufla.jpg}%

    \vspace{1cm}

    \Large\theauthor
    
    \vspace{\stretch{1}}

    \bfseries\LARGE\MakeUppercase\thetitle

    \vspace{\stretch{1}}

    \Large\MakeUppercase\olocal\\
    \MakeUppercase\thedate
    \null
    \cleardoublepage
\end{titlepage}
}

\newcommand \fazerfolhaderosto{
    \begin{titlepage}
        \newpage

        \begin{center}
            \MakeUppercase\theauthor

            \vspace{\stretch{1}}

            \textbf{\MakeUppercase\thetitle}

            \vspace{\stretch{1}}
            
            \hfill\begin{minipage}{8cm}{\otipo}\end{minipage} % tipo da monografia

            \vspace{\stretch{.58}}

            \ifthenelse{\equal{\@orientador}{\@empty}} % % se orientador nao existir,
            {}%                                            % nao produza nada
            {\orient\\
             Orientadora
            }

            \vspace{\stretch{.86}} 

            \textbf{\MakeUppercase\olocal}\\ 
            \textbf{\MakeUppercase\thedate}
            \end{center}
    \end{titlepage}
}

\newcolumntype{P}[1]{>{\centering\arraybackslash}p{#1}}
\graphicspath{{./figuras/}}

\newcommand{\tabitem}{~~\llap{\textbullet}~~}
\newcommand{\R}{$^{\tiny{\textregistered}}$~}


%==============================================================================

% Meus comandos

\newcommand{\citeay}[1]{\citeauthoronline{#1} (\citeyear{#1})}
\newcommand{\citeayp}[1]{(\citeauthoronline{#1} \citeyear{#1})}
\newcommand{\citean}[1]{\citeauthoronline{#1} \cite{#1}}

%==============================================================================
% Começar o documento

%\newsubfloat{figure}

%\includeonly{capitulos/materiais_metodos,capitulos/resultados}
\newcommand{\pa}[1]{\left(#1 \right)}
\begin{document}

\fazercapa

\fazerfolhaderosto

\tableofcontents

\clearpage

\pagestyle{ufla}

%==============================================================================
% incluindo os capitulos

\newcommand{\citeay}[1]{\citeauthoronline{#1} (\citeyear{#1})}
\newcommand{\citeayp}[1]{(\citeauthoronline{#1} \citeyear{#1})}
\newcommand{\citean}[1]{\citeauthoronline{#1} \cite{#1}}

\chapter{Introdução}\label{intro}

\section{Referencial Teórico}


Segundo \citeay{Munson04} ``a mecânica dos fluidos é a parte 
da mecânica aplicada que se dedica à
análise do comportamento dos líquidos e gases, tanto em equilíbrio quanto em 
movimento. Obviamente, o escopo da mecânica
dos fluidos abrange um vasto conjunto de 
problemas.''

Na análise do escoamento de fluidos é necessário o 
balanço de massa e de momento para desenvolvimento de equações
apropriadas aos fenômenos a serem analisados \citeayp{Cremasco14}.
A descrição matemática da fluidodinâmica é dada pelas 
equações de massa \eqref{eqcontinuidade}
e equações do movimento \eqref{eqQDM}.


\begin{equation}\label{eqcontinuidade}
    \frac{\partial\rho}{\partial t} + 
    \vec{\nabla} \cdot \rho \textbf{u}  = 0
\end{equation}

\begin{equation}\label{eqQDM}
    \rho \left(
        \frac{\partial \rho}{\partial t} + 
    \textbf{u} \cdot \vec{\nabla}\rho\textbf{u}
    \right)
    = 
    \rho\vec{\nabla}\varphi -
    \vec{\nabla} P +
    \vec{\nabla}\cdot \tau
\end{equation}

onde $\rho$ é a massa específica do fluido, \textbf{u} é a sua velocidade,
$\vec{\nabla}\varphi$ é o vetor intensidade do campo, $\vec{\nabla} P$ é o gradiente
de pressão e $\tau$ é a tensão extra, dinâmica ou viscosa.

Ao manipular as equações \eqref{eqcontinuidade} e 
\eqref{eqQDM}, considerando que o fluido se encontra em
regime permanente, é invíscido e incompressível, chega-se a
equação simplificada para a energia mecânica:

\begin{equation}\label{eqBern}
    \left(
        z + \frac{P}{\rho g} + \frac{u^2}{2g}
    \right)_{1}
    =
    \left(
        z + \frac{P}{\rho g} + \frac{u^2}{2g}
    \right)_{2} + h_L

\end{equation}

onde os subíndices $1$ e $2$ indicam dois pontos de escoamento dinstintos, e $h_L$ é
a perda de carga dada em unidades de comprimento.

Conforme \citeay{Cremasco14}, um fluido em um sistema de 
escoamento que apresente 
acessórios estará sujeito a uma perda de
carga, uma vez que, segundo o autor, quando o fluido entra em contato com 
um acessório ocorre a separação de uma camada
do escoamento, formando correntes turbulentas que dissipam energia mecânica. 
Essa energia dissipada representa a conversão irreversível de energia mecânica 
em energia térmica não desejada e a perda de energia por transferência de 
calor \citeayp{Fox14}.


A perda de carga oriunda de um acessório em um sistema de escoamento é definida por meio da equação:

\begin{equation}
    h_L = k_f \frac{u^2}{2g}
\end{equation}

onde $k_f$ é o coeficiente de perda de carga localizada característico do acessório, 
$u$ é a velocidade média do fluido e $g$ é a aceleração local da gravidade. 

Ainda conforme o autor, $k_f$ é constante apenas em regime turbulento. Podem ser 
encontradas na literatura tabelas que apresentam os valores de $k_f$ para válvulas
e acessórios específicos.

De acordo com \citeay{Potter15}, o regime turbulento é caracterizado 
pelo movimento
aleatório das partículas do fluido na
trajetória do escoamento. Tal regime pode ser previsto por um número 
adimensional denominado por número de Reynolds,
o qual é calculado, para tubulações cilíndricas, conforme:

\begin{equation}
    Re = \frac{u D}{\nu}
\end{equation}
 
onde $D$ é o diâmetro característico da tubulação e $\nu$ é a viscosidade 
cinemática do fluido.

Ainda conforme \citeay{Potter15}, valores para o número de 
Reynolds acima de 2000 são 
características de um escoamento em regime
turbulento, sendo este o valor mínimo usado na maioria das aplicações em engenharia.
Entretanto, tal valor pode variar em decorrência da geometria e da 
rugosidade das paredes da tubulação.

Portanto, além do cálculo adimensional, pode ser realizada uma análise 
qualitativa por meio do despejo de partículas
corantes na corrente de escoamento, avaliando se ocorre uma mistura imediata, ou 
não. Caso não, o regime é laminar.

\section{Aplicabilidade}

O cálculo da perda de carga é necessário em qualquer tipo de planejamento
e instalação de tubulações, uma vez que o mesmo está relacionado diretamente
com o matérial utilizado e as características da tubulação.

 Algumas situações que requerem uma análise da perda de carga podem ser 
exemplificadas como: possível cavitação em
bombas, danos à tubulação de edifícios, insuficiência na captação de água e/ou 
esgoto para tratamento e controle de
fluxo em reatores e outros 
equipamentos da indústria química.

\section{Objetivos}

O objetivo desse projeto é determinar a perda de carga e o coeficiente de 
perda de carga localizada em um tubo conectado
a dois reservatórios com um fluxo de água constante através de acessórios 
(individualmente) como registros e cotovelos a
partir da variação de pressão entre os pontos anterior e 
posterior aos acessórios.

% \chapter{Motivação}\label{motivacao}
Para o desenvolvimento de um projeto industrial que envolva transporte de 
fluidos é necessário o conhecimento do
estado energético-mecânico do fluido para que o mesmo possa ser deslocado 
conforme a necessidade do sistema.  Avaliando,
neste caso, a necessidade ou não da utilização de sistemas fluidomecânicos 
(bombas, compressores, por exemplo). 

A aferição do estado energético-mecânico do fluido em um determinado ponto de 
uma tubulação, por exemplo, requer a
investigação precisa dos pontos na tubulação por onde passou e o estado 
inicial do mesmo. Desta forma, elencar-se-á as
perdas de carga envolvidas durante o processo e, poder-se-á, então, determinar 
o estado de um fluido em um ponto e
avaliar se será necessário a injeção, ou remoção, de carga (energia) ao 
sistema. A presença de acessórios em uma
tubulação é quase que inevitável, seja por necessidades de otimização de 
espaço, seja por segurança ou por qualquer
outro motivo que seja necessária a instalação de uma válvula, por exemplo. 
Portanto, a investigação da perda de carga em
um escoamento é vital ao desenvolvimento de um projeto que atenda, com 
qualidade e segurança, o objetivo do sistema.

%%% Local Variables:
%%% mode: latex
%%% TeX-master: "../main_archive"
%%% End:
 % É para deixar?
\chapter{Descrição do Projeto}\label{descricao}

\section{Materiais Utilizados}

 \begin{table}[H]\label{t:custos}
     \centering
     \begin{tabular}{l c }
         \toprule
         \textbf{Material} & \textbf{Valor (R\$)} \\
         \midrule
         Adaptador (conexão entre balde e cano) & 1,00 - 5,00 \\
         Baldes & 20,00 - 30,00 \\
         Canos PVC & 25,00 - 50,00 \\
         Cola Durepoxi & 10,00 - 20,00 \\
         Cotovelos 90$^\circ$ & 1,00 - 5,00\\
         Cronômetro & 10,00 - 20,00 \\
         Registro rosqueável & 30,00 - 50,00 \\
         Suportes para a estrutura (madeira) & 20,00 - 50,00 \\
         Trena & 10,00 - 15,00 \\
         Tubos transparentes (manômetro) & 15,00 - 20,00 \\
         Válvula esférica & 30,00 - 50,00 \\
         \midrule
         \textbf{TOTAL} & 172,00 - 315,00\\
         \bottomrule
     \end{tabular}
     \caption{Custos}
 \end{table}

\section{Montagem}

A estrutura se inicia com um recipiente posicionado a uma elevação 
considerável do nível convencionado como $z = 0$. Esse
recipiente será um balde de alta capacidade volumétrica, o qual possuirá um ``ladrão''
na sua cota máxima, afim de evitar-se transbordamento do fluido excedente e também
estará ligado, em sua parte inferior, 
à tubulação de PVC através de um
adaptador, e após essa conexão será instalado um registro para controle da 
saída do fluido. A seguir, a tubulação fará
uma curva de $90^\circ$, descendo até o nível $z = 0$, 
onde fará uma nova curva de $90^\circ$
e prosseguirá reta até o local onde será
acoplado o acessório escolhido. Inicialmente, será analisada a perda de carga 
em uma válvula esférica e para tal é
necessário medir a diferença de pressão no tubo antes e depois da
válvula, com o auxílio de um manômetro.

Manômetros são instrumentos utilizados para determinação da diferença de 
pressão entre dois pontos de um fluido através
da diferença de elevação entre estes dois pontos. Dessa forma, o manômetro 
será constituído de dois tubos transparentes
acoplados ao tubo de PVC, um deles um pouco antes da válvula e o outro um 
pouco depois. Nesses tubos serão medidas as colunas de líquido. 

Em seguida, a estrutura irá terminar com a tubulação despejando o fluido em um 
outro recipiente de capacidade
volumétrica conhecida de forma que, através da cronometragem do tempo que esse 
recipiente leva para completar seu volume, seja possível 
determinar a vazão volumétrica. 



\chapter{Procedimento Experimental}
\label{chap:procedimento}

% Inicialmente, com a ajuda de um profissional, furou-se um balde de 20L em dois
% lugares: um localizado próximo à borda, para o ladrão, e outro próximo à base,
% para permitir a entrada de água na tubulação.

%% ===================================================

Inicialmente, antes das medições, o fluxo de água da torneira foi controlado de
modo que o regime de escoamento se tornasse estacionário, ou seja, não houvesse
variação no tempo. Após isso, foram feitas várias medições para cada abertura da
válvula esférica, sendo elas: totalmente aberta e uma abertura de 45 graus.
Essas aberturas foram calculadas com o auxílio de um transferidor, acoplado à
válvula. 

Para realização do experimento, foram realizadas, simultâneamente, a medição da
vazão e da diferença de pressão causada pela válvula. Para a vazão, foram
demarcados em um balde de 15L níveis correspondentes a volumes de 8 e 12 litros;
também foi coletado o tempo necessário para que a água que deixa a tubulação
preencha os respectivos volumes, com auxílio de cronômetros.  Com o intuito de
minimizar erros, foram utilizados dois cronômetros e, posterioremente,
calculou-se a média aritmética dos tempos obtidos. Já para aquisição dos dados
de diferença de pressão, mediu-se, com trenas, as alturas de coluna de água em
cada uma das mangueiras.

O procedimento supracitado foi repetido cinco vezes, para cada abertura da
válvula, afim de garantir maior precisão nos resultados. E os dados obtidos
serão tratados em (REFERENCIAR)


%%% Local Variables:
%%% mode: latex
%%% TeX-master: "../main_archive"
%%% End:

\chapter{Resultados e Discussão}
\label{chap:resultados}

\section{Cálculo da vazão}
\label{sec:vazao}

Conforme procedimento experimental, apresenta-se a tabela \ref{t:tempos}, a
qual relaciona para cada amostragem os tempos de preenchimento dos volumes
analisados pela razão volumétrica do fluido de trabalho -- água.

\begin{table}[H]
\centering
\caption{Tempos amostrais de preenchimento.}
\label{t:tempos}
\begin{tabular}{|c|c|c|c|c|}
\hline
\multicolumn{5}{|c|}{\textbf{Tempo para preenchimento (s)}}                                      \\ \hline
\multicolumn{1}{|l|}{} & \multicolumn{2}{c|}{\textbf{0$^\circ$}}    & \multicolumn{2}{c|}{\textbf{45$^\circ$}}   \\ \hline
\textbf{Amostra}       & \textbf{8L} & \textbf{12L} & \textbf{8L} & \textbf{12L} \\ \hline
1                      & 38,780          & 55,500           & 84,455          & 126,570          \\ \hline
2                      & 47,570          & 71,000           & 85,290          & 127,245          \\ \hline
3                      & 46,140          & 70,520           & 85,980          & 127,050          \\ \hline
4                      & 47,185          & 71,155           & 85,175          & 127,205          \\ \hline
\end{tabular}
\end{table}

As vazões obtidas para cada amostragem foram calculadas considerando os tempos
necessários para que o volume ocupado pelo fluido no recipiente refratário
utilizado alcançasse 8 e 12L, utilizando a seguinte equação:

\newcommand{\Vol}{\mathop{\ooalign{\hfil$V$\hfil\cr\kern0.08em--\hfil\cr}}\nolimits}
\begin{equation}
  Q = \frac{\Vol}{\Delta t}
\end{equation}

onde $Q$ é a vazão e $\Vol$ é o volume preenchido durante o intervalo de tempo
$\Delta t$. Assim, utilizando-se dos dados constantes na tabela \ref{t:tempos},
apresenta-se a tabela seguinte, a qual relaciona a vazão do sistema para cada
amostra.

\begin{table}[H]
\centering
\caption{Vazão amostral por ângulo de abertura.}
\label{t:vazao}
\begin{tabular}{|c|c|c|c|c|c|c|}
\hline
\multicolumn{7}{|c|}{\textbf{Vazão ($m^3/s$)}}                                                               \\ \hline
\textbf{}        & \multicolumn{3}{c|}{\textbf{$0^\circ$}}     & \multicolumn{3}{c|}{\textbf{$45^\circ$}}    \\ \hline
\textbf{Amostra} & \textbf{8L} & \textbf{12L} & \textbf{Média} & \textbf{8L} & \textbf{12L} & \textbf{Média} \\ \hline
1                & 2,06292E-04 & 2,16216E-04  & 2,11254E-04    & 9,47250E-05 & 9,48092E-05  & 9,47671E-05    \\ \hline
2                & 1,68173E-04 & 1,69014E-04  & 1,68594E-04    & 9,37976E-05 & 9,43063E-05  & 9,40519E-05    \\ \hline
3                & 1,73385E-04 & 1,70164E-04  & 1,71775E-04    & 9,30449E-05 & 9,44510E-05  & 9,37479E-05    \\ \hline
4                & 1,69545E-04 & 1,68646E-04  & 1,69090E-04    & 9,39243E-05 & 9,43359E-04  & 9,41301E-05    \\ \hline
5                & 1,68315E-04 & 1,69803E-04  & 1,69059E-04    & 9,48148E-05 & 9,42840E-04  & 9,45494E-04    \\ \hline
\end{tabular}
\end{table}

\section{Cálculo da velocidade do fluido}
\label{sec:VelFluido}

A partir dos dados da tabela \ref{t:vazao}, calculou-se as respectivas
velocidades do fluido no sistema para cada amostra utilizando a equação

\begin{equation}
  v = \frac{4 Q}{\pi D^2}
\end{equation}

onde $v$ é a velocidade média do fluido e $D$ é o diâmetro hidráulico da
tubulação.

Tais velocidades foram consideradas constantes em todo o sistema, pois o fluido
de trabalho foi a água no estado líquido, portanto incompressível e a área da
seção transversal do escoamento interno manteve-se constante (tubo rígido). Tais
velocidades estão relacionadas às respectivas amostragens, conforme a seguinte
tabela.

\begin{table}[H]
\centering
\caption{Velocidade e número de Reynolds por ângulo de abertura.}
\label{t:Reynolds}
\begin{tabular}{|c|c|c|c|c|}
\hline
\multicolumn{5}{|c|}{\textbf{Velocidade e  número de Reynolds}}                                                      \\ \hline
\textbf{}        & \multicolumn{2}{c|}{\textbf{$0^\circ$}}         & \multicolumn{2}{c|}{\textbf{$45^\circ$}}        \\ \hline
\textbf{Amostra} & \textbf{Velocidade ($m/s$)} & \textbf{Reynolds} & \textbf{Velocidade ($m/s$)} & \textbf{Reynolds} \\ \hline
1                & 0,67244256                  & 13397,8475        & 0,30165304                  & 6010,18103        \\ \hline
2                & 0,53665026                  & 10692,3012        & 0,29937664                  & 5964,82561        \\ \hline
3                & 0,54677655                  & 10894,0590        & 0,29840899                  & 5945,54597        \\ \hline
4                & 0,53824820                  & 10724,1389        & 0,29962540                  & 5969,78182        \\ \hline
5                & 0,53813161                  & 10721,8157        & 0,30096016                  & 5996,37588        \\ \hline
\end{tabular}
\end{table}

\section{Cálculo do número de Reynolds (Re)}
\label{sec:Reynolds}

O $Re$ do sistema foi calculado para cada amostragem considerando a equação
\eqref{numReynolds}, na qual a variável geométrica característica para tubos
circulares é o diâmetro. As propriedades do fluido foram consultadas em
\citeay{Geankoplis93} considerando a pressão local igual a 1 atm ($101,325$ kPa) e
a $20^\circ$ C ($293,15$ K). Tais valores são apresentados na tabela
\ref{t:Reynolds}.

Conforme a tabela \ref{t:Reynolds}, cada amostragem apresentou $Re > 2000$ (valor
de referência utilizado conforme descrito no referencial teórico) caracterizando
para cada amostragem escoamento em regime turbulento.


\section{Cálculo do coeficiente de perda localizada ($k_f$)}
\label{sec:carga}

Conforme \citeay{Cremasco14}, em regime turbulento o coeficiente de perda de
carga localizada ($k_f$) da equação \eqref{PerdaDeCarga} é constante. Portanto,
tal coeficiente pode ser avaliado ao ser isolado da equação
\eqref{PerdaDeCarga}, uma vez que a perda de carga total referente à
configuração do sistema avaliado pode ser determinada e igualada a perda de
carga localizada, pois as perdas de cargas maiores foram desprezadas em
decorrência da pequena distância entre os pontos analisados e da rugosidade
relativa do tubo (de PVC) ser desprezível (tubo liso).

Apresentam-se, então, as seguintes tabelas, cujas quais relacionam a amostragem
às respectivas alturas de pressão (Tabela \ref{t:alturas}) e os coeficientes
de perda de carga localizados calculados (Tabela \ref{t:kf}). Para cálculo dos
valores de perda de carga e do coeficiente de perda de carga, as seguintes
equações foram utilizadas:

\begin{equation}\label{eq:hL}
h_L = z_1 - z_2
\end{equation}

\begin{equation}\label{eq:kf}
  k_f = \frac{2 g h_L}{v^2}
\end{equation}

onde $z_1$ e $z_2$ são a altura da coluna de água coletada nos manômetros
situados antes e após a válvula esférica, respectivamente. A equação
\eqref{eq:hL} pode ser obtida através da equação \eqref{eqBern}, aplicando-a
corretamente para o sistema apresentado.

\begin{table}[H]
\centering
\caption{Altura manométrica amostral por ângulo de abertura.}
\label{t:alturas}
\begin{tabular}{|c|c|c|c|c|}
\hline
\multicolumn{5}{|c|}{\textbf{Alturas (cm)}}                                                           \\ \hline
\textbf{}        & \multicolumn{2}{c|}{\textbf{$0^\circ$}} & \multicolumn{2}{c|}{\textbf{$45^\circ$}} \\ \hline
\textbf{Amostra} & \textbf{z1 (cm)}   & \textbf{z2 (cm)}   & \textbf{z1 (cm)}    & \textbf{z2 (cm)}   \\ \hline
1                & 9,0                & 3,5                & 71,6                & 3,3                \\ \hline
2                & 7,6                & 3,6                & 72,2                & 3,1                \\ \hline
3                & 7,4                & 3,7                & 72,3                & 3,4                \\ \hline
4                & 7,5                & 4,0                & 72,5                & 3,6                \\ \hline
5                & 7,8                & 4,2                & 71,9                & 3,5                \\ \hline
\end{tabular}
\end{table}

\begin{table}[H]
\centering
\caption{Perda de carga e coeficiente de perda de carga amostral por ângulo de abertura.}
\label{t:kf}
\begin{tabular}{|c|c|c|c|c|}
\hline
\multicolumn{5}{|c|}{\textbf{Perda de Carga e Coeficiente de Perda de Carga}}                               \\ \hline
\textbf{}              & \multicolumn{2}{c|}{\textbf{$0^\circ$}} & \multicolumn{2}{c|}{\textbf{$45^\circ$}} \\ \hline
\textbf{Amostra}       & \textbf{$h_L$ (m)}   & \textbf{$k_f$}   & \textbf{$h_L$ (m)}    & \textbf{$k_f$}   \\ \hline
1                      & 0,055                & 2,38644435       & 0,683                 & 147,266607       \\ \hline
2                      & 0,040                & 2,72506149       & 0,691                 & 151,265967       \\ \hline
3                      & 0,037                & 2,42818046       & 0,689                 & 151,807917       \\ \hline
4                      & 0,035                & 2,37029210       & 0,689                 & 150,577813       \\ \hline
5                      & 0,036                & 2,43907135       & 0,684                 & 148,162087       \\ \hline
\textbf{Desvio Padrão} & 0,00826438           & 0,14550331       & 0,00349285            & 1,99282035       \\ \hline
\textbf{Média}         & 0,0406               & 2,52190594       & 0,6872                & 149,812607       \\ \hline
\end{tabular}
\end{table}

A NBR 12214/1992 (Projeto de sistema de bombeamento de água para abastecimento
público, \cite{NBR12214}), em seu anexo B, traz valores para os coeficientes de
perda de carga localizada para a equação \eqref{PerdaDeCarga} para o acessório
utilizado -- válvula esférica. Para a abertura de $0^\circ$ (totalmente aberta),
$k_f$ = 0. Já para abertura de $45^\circ$, $k_f$ = $31,2$.

Conforme a tabela acima, para a abertura de $0^\circ$ o $k_f$ obtido foi igual a
$2,5219 \pm 0,1455$. Já para a abertura de $45^\circ$, o $k_f$ obtido foi igual
a $149,8126 \pm 1,9928$. Ambos os resultados diferiram, estatisticamente ao
nível de 95\% de confiança, dos respectivos valores referenciados pela norma
supracitada.

Tais diferenças podem ser justificadas devido às adaptações realizadas no
sistema analisado, que podem ter prejudicado os resultados para as duas
configurações. Além disso, há que se considerar os erros sistemáticos e
aleatórios presentes durante a aquisição dos dados.

Dos erros sistemáticos, citam-se:
\begin{itemize}
\item Instrumental: Na determinação do volume do refratário inferior, o qual
  pode ter sido mal calibrado;
\item Observacional: Erro paralaxe na leitura dos tempos;
\item Teórico: Oriundo da simplificação das equações que modelam o sistema.
\end{itemize}

Dos erros aleatórios, cita-se:
\begin{itemize}
\item Observacional: Erro na leitura da altura da coluna de fluido.
\end{itemize}

%%% Local Variables:
%%% mode: latex
%%% TeX-master: "../main_archive"
%%% End:

% \chapter{Discussão}

A ESCREVER

%%% Local Variables:
%%% mode: latex
%%% TeX-master: "../main_archive"
%%% End:
  % 
\chapter{Conclusão}
\label{chap:conclusao}

Diante do exposto e da análise dos resultados obtidos supracitada, conclui-se
que a determinação experimental do coeficiente de perda de carga localizada
requer condições laboratoriais adequadas de modo que as adaptações necessárias
sejam mínimas a fim de maximizar a acurácia dos resultados. No que tange o
propósito didático à compreensão dos mecanismos físicos da mecânica dos fluidos,
avalia-se de modo satisfatório, uma vez que a execução do presente projeto
proporcionou a aplicabilidade da teoria, como também fomentou o raciocínio à
solução de problemas reais de um engenheiro.

%%% Local Variables:
%%% mode: latex
%%% TeX-master: "../main_archive"
%%% End:


%==============================================================================

% Incluindo bibliografia
%\bibliographystyle{plain}             % estilo para labels em numeros
%\bibliographystyle{alpha}             % estilo para labels em iniciais
\bibliographystyle{abntex2-num}         %% estilo para referências usando ABNT, 
 %para usar o abntex2-alf precisa arrumar o nome dos autores na bibliografia
%\bibliographystyle{abntex2-alf}           % estilo para referências usando ABNT, 

%inclui Referências Bibliográficas
\referencias
\bibliography{bibliografia/b}

%==============================================================================
% Fim do texto
\end{document}

%%% Local Variables:
%%% mode: latex
%%% TeX-master: t
%%% End:
